% ============================================================================
%  Section VII — Discussion and Limitations  (IEEE TIM, IEEEtran)
%  Passive voice / third person throughout. Interpretation of the Section VI
%  results, practical guidance, and limitations. Requires: none beyond base.
% ============================================================================

\section{Discussion and Limitations}
\label{sec:discussion}

\subsection{Discussion}

\subsubsection{Evaluation methodology}
The optimism gap of Section~\ref{subsec:gap} is the central methodological
finding: for low-rate smartphone-based machine diagnosis, accuracies obtained under
random splitting of overlapping windows overstate deployable performance by tens
of accuracy points. The effect mirrors that documented for industrial-grade
vibration data~\cite{wheat2024impact, hendriks2022towards} and the broader
machine-learning leakage literature~\cite{kapoor2023leakage}, but its magnitude
here is large because each operating condition is represented by a single
continuous recording. A practical consequence is that the near-perfect accuracies
frequently reported for this class of data should be interpreted with caution; the
recording-wise and cross-condition figures obtained here provide a more
deployment-honest reference against which future methods can be compared.

\subsubsection{Ultra-low-bandwidth sensing as an asset}
The observation that sampling rates at or below $25$~Hz match or exceed the
accuracy obtained at $100$~Hz reframes the low sampling rate of consumer MEMS
sensors from a limitation into an advantage. The information required to separate
the coarse fault states resides in the low-frequency band identified in
Section~\ref{sec:charact}, so aggressive decimation reduces the telemetry and, by
extension, the energy of a wireless node without sacrificing accuracy; it can even
improve generalisation by discarding high-frequency content that the classifier
would otherwise overfit. The resource-optimal configuration established in
Section~\ref{subsec:tradeoff}---$25$~Hz sampling, the three raw channels, and a
$4$~s window---is therefore recommended as a default operating point for this
class of node.

\subsubsection{Model complexity and edge feasibility}
Under leakage-free evaluation the lower-capacity models generalised at least as
well as the Random Forest while occupying a kilobyte-scale footprint, so on-node
inference is practical on inexpensive microcontrollers. Combined with the link
budget, this yields per-link deployment guidance: Bluetooth Low Energy, Zigbee,
and NB-IoT nodes can stream raw or feature data to a gateway, whereas LoRaWAN
nodes, constrained by the duty cycle to the decision stream, must perform
classification locally. The alignment between the models that deploy most easily
and those that generalise best is convenient for practical systems.

\subsubsection{Sensor-level guidance}
The channel analysis of Section~\ref{sec:charact} carries a concrete
recommendation for smartphone-based acquisition: the raw acceleration channels
should be preferred over the gravity-compensated channels, because the on-device
gravity compensation---valuable for motion-tracking applications---attenuates
low-frequency content that is diagnostically useful for condition monitoring.

\subsubsection{Scope of applicability}
Taken together, the results delineate where low-cost, low-rate MEMS sensing is
appropriate. Coarse health states, such as a healthy machine, a broken rotor bar,
or a supply-voltage imbalance, are separable at $100$~Hz and below; the staging of
incipient bearing degradation, whose signatures lie in the high-frequency band, is
not, and remains the province of higher-rate, industrial-grade instrumentation.

\subsection{Limitations}
\label{subsec:limits}
Several limitations bound the scope of these findings. The dataset comprises a
single $1.1$~kW machine with laboratory-induced, discrete-severity faults, and the
two lubrication states were graded qualitatively; the reported accuracies are
therefore a deployment-honest baseline rather than a claim of field performance.
The small number of recordings ($36$) produces wide confidence intervals under
the recording-holdout protocols; although five-seed repetition stabilises the
central estimates, larger multi-machine datasets are required to narrow them. The
edge-node analysis quantifies communication load and model footprint but reports
energy only indirectly, since direct on-device measurement was outside the present
scope; a cycle-accurate estimation of inference latency and energy on a
representative microcontroller is left as future work. Finally, generalisation
across machines of different ratings and manufacturers was not assessed and
constitutes the most important direction for future validation.
